% ============================================================
% Do Morphological Auxiliary Losses Actually Improve Crack Segmentation?
% A Systematic Multi-Seed Study
% Target: Automation in Construction (elsarticle-num)
% Double-blind review version
% ============================================================

\documentclass[preprint,12pt]{elsarticle}

\usepackage{amsmath}
\usepackage{booktabs}
\usepackage{graphicx}
\usepackage{subcaption}
\usepackage{hyperref}
\usepackage{xcolor}
\usepackage{multirow}

\hypersetup{colorlinks=true, linkcolor=blue, citecolor=blue, urlcolor=blue}

\journal{Automation in Construction}

\begin{document}

\begin{frontmatter}

\title{Do Morphological Auxiliary Losses Actually Improve Crack Segmentation? A Systematic Multi-Seed Study}

\author[anon]{Anonymous}
\address[anon]{Removed for double-blind review}

\begin{abstract}
Auxiliary morphological supervision---skeleton distance transforms, topology-preserving losses, keypoint heatmaps---is widely adopted to improve crack segmentation in concrete infrastructure inspection, yet its effectiveness has not been rigorously validated with multi-seed experiments. We present a systematic evaluation of five auxiliary supervision strategies built as a progressive baseline ladder on SegFormer-B2 with a shared Feature Pyramid Network, three post-hoc graph extraction methods, direct neural graph prediction, and three standard domain generalisation methods (CORAL, DANN, MixStyle), all evaluated on concrete dam inspection imagery. Our central finding is negative: no auxiliary morphological supervision significantly improves segmentation over the mask-only baseline under multi-seed validation, and no domain generalisation method reduces the cross-site performance gap in a pseudo-domain setting. More critically, we identify three cases where single-seed evaluation produced misleading conclusions that were overturned by multi-seed validation: skeleton distance transform supervision appeared to gain +1.0\% mIoU (actual: +0.5\%$\pm$0.3\%, $p{=}0.145$), clDice appeared to lose $-$1.6\% (actual: +0.2\%$\pm$0.9\%, $p{=}0.822$), and MixStyle appeared to improve out-of-distribution performance by +52\% (actual: ${\sim}$0\%, $p{=}0.878$). These results demonstrate the fragility of single-seed comparisons and underscore the need for multi-seed validation as a minimum standard in this domain.
\end{abstract}

\begin{keyword}
crack segmentation \sep auxiliary supervision \sep multi-seed validation \sep domain generalisation \sep graph extraction \sep concrete inspection
\end{keyword}

\end{frontmatter}


% ============================================================
% INTRODUCTION (preserved verbatim from intro_relwork.tex)
% ============================================================

\section{Introduction}
\label{sec:introduction}

Automated detection and quantification of surface defects---primarily cracks and spalling---is a critical task in structural health monitoring of concrete infrastructure~\cite{koch2015review, mosalam2024dlshm}. Manual visual inspection remains labour-intensive, subjective, and difficult to scale across geographically distributed assets such as dams, bridges, and tunnels. Over the past decade, deep learning has transformed this domain: convolutional and transformer-based semantic segmentation models now routinely achieve high per-pixel accuracy on benchmark crack datasets~\cite{kulkarni2022crackseg9k, dorafshan2018sdnet, zhang2025crackreview}.

The dominant paradigm employs encoder--decoder architectures---U-Net~\cite{ronneberger2015unet}, DeepLabV3+~\cite{chen2018deeplabv3plus}, and more recently SegFormer~\cite{xie2021segformer}---trained with standard pixel-level losses such as cross-entropy and Dice. While these models perform well on in-distribution data, two persistent challenges motivate ongoing research. First, cracks are thin, elongated structures whose \emph{topological} integrity (connectivity, branching) is poorly captured by area-based loss functions~\cite{shit2021cldice, hu2019topology}. This has led to a growing body of work on topology-preserving auxiliary losses and morphological supervision heads that aim to improve segmentation quality beyond what per-pixel objectives can achieve. Second, models trained on imagery from one inspection site often generalise poorly to other sites, sensors, or surface conditions---a cross-site domain gap that remains largely unaddressed in the structural inspection literature~\cite{bukhsh2021damage, hassani2025dashm}.

To tackle the first challenge, researchers have proposed auxiliary losses and supervision signals drawn from crack morphology: centerline Dice (clDice)~\cite{shit2021cldice}, persistent-homology-based topological losses~\cite{clough2020topological, hu2019topology}, skeleton distance transform regression~\cite{audebert2019distance}, and keypoint heatmaps for endpoints and junctions~\cite{sun2019crackjunction}. These methods are widely assumed to improve segmentation quality, and published results often show gains of one to several percentage points in mean Intersection-over-Union (mIoU) or similar metrics. However, a critical methodological concern undermines confidence in these findings: the vast majority of studies report results from a \emph{single random seed}, making it impossible to distinguish genuine methodological improvements from seed-dependent variance~\cite{picard2021seed, henderson2018rl}. Recent work has demonstrated that random seed selection alone can produce performance differences comparable in magnitude to claimed improvements, both in computer vision generally~\cite{picard2021seed} and in civil engineering applications specifically~\cite{alsabbagh2024randomness}.

For the second challenge, standard domain generalisation (DG) methods---including CORAL~\cite{sun2016coral}, DANN~\cite{ganin2016dann}, and MixStyle~\cite{zhou2021mixstyle}---have been developed to learn domain-invariant representations. While these methods have demonstrated effectiveness in multi-source settings (e.g., photo/art/sketch domains), their applicability to single-source infrastructure inspection, where pseudo-domains must be constructed from within-dataset partitions, has not been systematically evaluated.

This paper presents a rigorous, multi-seed empirical study that challenges both lines of assumption. We construct a progressive baseline ladder on SegFormer-B2~\cite{xie2021segformer} with a shared Feature Pyramid Network~\cite{lin2017fpn}, systematically adding morphological supervision---clDice, skeleton distance transform regression, keypoint heatmaps, and width estimation---and evaluating each configuration with pre-registered go/no-go gates. We further evaluate three standard DG methods under a pseudo-domain setting using difficulty-based data partitioning. All experiments are validated with three independent random seeds, and statistical significance is assessed via paired $t$-tests.

Our central finding is \textbf{negative}: no form of auxiliary morphological supervision significantly improves crack segmentation over the mask-only baseline when validated with multiple seeds, and no standard DG method reduces the cross-site domain gap under pseudo-domain single-source training. More critically, we identify \textbf{three cases of misleading single-seed results} that were overturned by multi-seed validation:
\begin{enumerate}
    \item Skeleton distance transform supervision appeared to improve mIoU\textsubscript{fg} by +1.0\% in a single seed, but multi-seed evaluation revealed a non-significant +0.5\%$\pm$0.3\% ($p=0.145$).
    \item clDice appeared to degrade performance by $-$1.6\% in a single seed, but multi-seed evaluation showed a non-significant +0.2\%$\pm$0.9\% ($p=0.822$).
    \item MixStyle appeared to improve out-of-distribution performance by +52\% in a single seed, but multi-seed evaluation showed statistical equivalence to the ERM baseline ($p=0.878$).
\end{enumerate}

These results carry implications beyond our specific experimental setting. They suggest that the crack segmentation literature may contain a non-trivial number of fragile claims built on single-seed comparisons, and they underscore the need for multi-seed validation as a minimum standard of evidence for claimed improvements.

The contributions of this work are as follows:
\begin{itemize}
    \item A systematic evaluation of five auxiliary morphological supervision strategies (clDice, skeleton distance transform, endpoint/junction detection, width regression, and their combination) for concrete defect segmentation, demonstrating that none significantly improves over a mask-only baseline under multi-seed validation.
    \item A multi-seed evaluation of three standard domain generalisation methods (CORAL, DANN, MixStyle) under pseudo-domain single-source training, showing that all are statistically equivalent to empirical risk minimisation.
    \item Identification of three concrete cases where single-seed evaluation produced misleading conclusions, providing an empirical demonstration of the fragility of single-seed comparisons in this domain.
    \item A multi-tier graph evaluation protocol for assessing the quality of crack graphs extracted from segmentation predictions, with strict, relaxed, and lenient tolerance tiers.
\end{itemize}


% ============================================================
% RELATED WORK (preserved verbatim from intro_relwork.tex)
% ============================================================

\section{Related Work}
\label{sec:related_work}

We review three research areas that our systematic study directly evaluates: topology-preserving losses for crack segmentation, crack graph and skeleton extraction methods, and domain generalisation for infrastructure defect detection.

\subsection{Topology-Preserving Losses for Crack Segmentation}
\label{sec:rw_topology}

Standard segmentation losses such as cross-entropy and Dice operate on a per-pixel basis and are agnostic to the topological structure of predicted regions. For thin, elongated objects like cracks, vessels, and roads, this can result in segmentations that achieve high volumetric overlap but contain topologically incorrect disconnections or spurious connections. This limitation has motivated a substantial body of work on topology-aware loss functions.

Hu et al.~\cite{hu2019topology} introduced a topological loss based on persistent homology that penalises differences in Betti numbers between predicted and ground-truth segmentations, with demonstrated improvements on neuron membrane and vessel segmentation. Clough et al.~\cite{clough2020topological} extended this approach with a differentiable persistent homology loss applicable to both 2D and 3D segmentation tasks. While theoretically principled, persistent-homology-based losses are computationally expensive and can be difficult to combine with standard losses during training~\cite{hu2021topology}.

An alternative line of work exploits morphological skeletons rather than algebraic topology. Shit et al.~\cite{shit2021cldice} proposed the centerline Dice (clDice) loss, which computes Dice overlap on soft morphological skeletons of the prediction and ground truth. The clDice loss is computationally efficient and theoretically guarantees topology preservation up to homotopy equivalence for binary segmentation. It has since been widely adopted for vessel, road, and crack segmentation. Kirchhoff et al.~\cite{kirchhoff2024skeleton} proposed the Skeleton Recall Loss, which performs tubed skeletonisation on the ground truth and computes a soft recall loss, reducing computational overhead by over 90\% compared to clDice while achieving competitive or superior connectivity preservation.

In the crack detection domain specifically, He and Li~\cite{he2024beyond} evaluated skeleton-based losses for preserving connectivity in crack segmentation, demonstrating improvements on the EdmCrack600 and CrackForest datasets using U-Net++ and SegFormer backbones. Jing et al.~\cite{jing2025topology} proposed TopoM-CrackNet, which combines persistent homology with a Mamba-enhanced U-Net for pavement crack detection with topology preservation, published in \textit{Automation in Construction}. Zhang and Qi~\cite{zhang2024irfusionformer} introduced a topology-based loss within an RGB-thermal fusion framework for pavement crack segmentation. Xu et al.~\cite{xu2021cploss} proposed CP-loss for connectivity preservation in road curb detection, which focuses on regions of disconnectivity during optimisation.

A critical observation across this literature is that the vast majority of studies report results from a single random seed. While the methods are often evaluated on multiple datasets, the lack of multi-seed validation within each dataset makes it difficult to assess whether reported improvements exceed the variance attributable to random initialisation. Our study directly addresses this gap by evaluating clDice~\cite{shit2021cldice} under multi-seed conditions and finding that its effect on crack segmentation is statistically indistinguishable from zero.

\subsection{Crack Graph and Skeleton Extraction}
\label{sec:rw_graph}

Beyond pixel-level segmentation, practical crack assessment requires measurable morphological features: centerline geometry, branching topology, junction locations, crack width, and crack length~\cite{li2024crackquantification, lin2025crackmorph}. Extracting a graph representation---where nodes represent endpoints and junctions, and edges represent crack segments---from a segmentation mask is a natural approach to obtaining such features.

Traditional approaches rely on morphological skeletonisation of binary crack masks, followed by graph extraction from the resulting skeleton pixels~\cite{wang2025crackskeletonization}. Endpoints (degree-1 pixels) and junctions (degree$\geq$3 pixels) are identified directly from the skeleton topology, and connected skeleton segments between these keypoints form graph edges. This pipeline is straightforward but sensitive to segmentation noise: small errors in the predicted mask can produce dramatically different skeletons, leading to fragile graph representations. Li and Yuan~\cite{li2025cracktopology} recently demonstrated a graph-based quantification framework for automated topological analysis of crack networks in road maintenance.

Distance transform (DT) maps provide an alternative representation that encodes the distance of each pixel to the nearest boundary~\cite{audebert2019distance}. Ridge detection on DT maps can recover crack centerlines without explicit skeletonisation, and local maxima correspond to the widest points of crack segments. Audebert et al.~\cite{audebert2019distance} showed that auxiliary DT regression can spatially regularise semantic segmentation predictions, and this approach has been adopted in various infrastructure inspection contexts.

Junction detection in crack images has been addressed through correlation structure analysis and tensor voting~\cite{sun2019crackjunction}, achieving high correctness and completeness on pavement images. More recently, Liu and Li~\cite{liu2024ltggnet} proposed LTGGNet, which uses local topology and global grouping attention for concrete crack detection. Lin and Zheng~\cite{lin2025crackmorph} introduced CrackMorph-XAI-Net, an explainable framework that integrates topology-preserving skeleton extraction with junction detection via Gaussian heatmap regression, producing measurable morphological descriptors. Kuttner and Grunwald~\cite{kuttner2025cracktopo} proposed extensions to the crack topology score for evaluating the topological quality of crack segmentations.

A common assumption in this literature is that auxiliary supervision targeting morphological features---skeleton DT regression, keypoint detection---should improve the quality of downstream graph extraction. Our study tests this assumption empirically: we find that skeleton DT supervision (B2) does not significantly improve post-hoc graph extraction quality compared to a mask-only baseline (B0), with all Wilcoxon signed-rank tests yielding $p > 0.23$. We further find that direct neural graph prediction using CenterNet-style~\cite{zhou2019centernet} node heatmaps and edge classifiers, while architecturally sound (demonstrated by overfit gate passage), is limited by the quality of automatically generated training labels.

\subsection{Domain Generalisation for Infrastructure Defect Detection}
\label{sec:rw_dg}

Deep learning models for structural inspection are typically trained on data from a single site or a small number of sites and are expected to generalise to unseen infrastructure. In practice, substantial performance degradation is observed when models are deployed on imagery from different sites, sensors, or environmental conditions~\cite{koch2015review, bukhsh2021damage}. This cross-site domain gap is well documented: Zhang and Wei~\cite{zhang2025crackreview} note that generalisability across datasets remains one of the most significant open challenges in crack detection. Hassani and Dackermann~\cite{hassani2025dashm} provide a systematic review of domain adaptation methods for structural health monitoring, covering both vibration-based and vision-based approaches.

Domain adaptation (DA) methods, which assume access to unlabelled target data during training, have been applied to crack segmentation with some success. Jiang and Yuan~\cite{jiang2025udacrack} proposed a cross-domain stylisation framework with dual adversarial feature learning for crack segmentation across different surface types. Jiang and Deng~\cite{jiang2025crossdataset} demonstrated cross-dataset semantic segmentation for composite crack detection using unsupervised transfer learning. Bukhsh et al.~\cite{bukhsh2021damage} evaluated in-domain and cross-domain transfer learning for structural damage detection, finding that cross-domain transfer is substantially more challenging.

Domain generalisation (DG) methods, which aim to learn domain-invariant representations \emph{without} access to target data, represent a more challenging but practically relevant setting. Three widely used DG methods are: (1)~Deep CORAL~\cite{sun2016coral}, which aligns second-order feature statistics across domains; (2)~DANN~\cite{ganin2016dann}, which uses a gradient reversal layer to learn features that are discriminative for the task but invariant to domain identity; and (3)~MixStyle~\cite{zhou2021mixstyle, zhou2023mixstyle}, which synthesises new domains by mixing instance-level feature statistics during training. These methods have been developed and validated primarily in multi-source DG benchmarks (e.g., PACS, Office-Home) where distinct visual domains are available.

A key practical constraint in infrastructure inspection is that labelled data often comes from a single source domain. In this single-source DG setting, pseudo-domains must be constructed---for example, by partitioning data based on difficulty, augmentation style, or other heuristics~\cite{elghazy2025pseudodg}. However, the effectiveness of pseudo-domain construction is not guaranteed: if the constructed pseudo-domains do not capture genuine distributional variation, DG methods may fail to learn meaningful domain-invariant features. Elghazy and Alaa~\cite{elghazy2025pseudodg} note that if a pseudo-domain's distribution is semantically irrelevant, it can disrupt the learning process and degrade model performance.

Our study evaluates all three DG methods (CORAL, DANN, MixStyle) in a single-source pseudo-domain setting where DamSegment training images are partitioned into difficulty-based tiers (Easy/Medium/Hard). Multi-seed evaluation with three seeds reveals that all three methods are statistically equivalent to the ERM baseline ($p=0.878$ for MixStyle vs.\ ERM on out-of-distribution mIoU\textsubscript{fg}), with a stable domain gap of approximately 55\%. This finding suggests that difficulty-based pseudo-domains do not capture the distributional characteristics needed for effective domain generalisation in concrete defect detection, and that bridging real cross-site gaps likely requires genuine multi-source data collection rather than algorithmic solutions applied to single-source data.


% ============================================================
% SECTION 3: METHODOLOGY
% ============================================================

\section{Methodology}
\label{sec:methodology}

This section describes the baseline architecture, the progressive auxiliary supervision ladder, the post-hoc graph extraction methods, and the domain generalisation methods evaluated in this study. An overview of the progressive baseline ladder is shown in Figure~\ref{fig:fig_baseline_ladder_overview}.

\begin{figure}[t]
\centering
\includegraphics[width=\columnwidth]{figures/fig_baseline_ladder_overview.png}
\caption{Overview of the progressive baseline ladder architecture. Starting from a SegFormer-B2 encoder with SharedFPN, auxiliary morphological heads are incrementally added: B0 (mask-only), B1a (+clDice loss), B2 (+skeleton distance transform head), B3 (+endpoint/junction heads), and B5 (+width head). Below, three post-hoc graph extraction methods operate on predicted masks and DT maps to produce crack graphs.}
\label{fig:fig_baseline_ladder_overview}
\end{figure}

\subsection{Baseline Architecture}
\label{sec:baseline_arch}

All experiments are built on a SegFormer-B2~\cite{xie2021segformer} backbone, which consists of a Mix Transformer (MiT) hierarchical encoder and a lightweight MLP decoder. The MiT encoder produces multi-scale feature maps at strides of 4, 8, 16, and 32, which the MLP decoder fuses into a unified feature representation. We augment this architecture with a shared Feature Pyramid Network (SharedFPN)~\cite{lin2017fpn} that produces multi-scale feature maps at the input resolution of $512 \times 512$ pixels, enabling flexible attachment of multiple task-specific heads.

The base configuration (B0) performs three-class semantic segmentation (background, crack, spalling) using a segmentation decoder head. The training objective combines a weighted cross-entropy loss with class weights $[0.2, 2.0, 3.0]$ to address severe class imbalance (background 96.6\%, crack 2.2\%, spalling 1.2\%) and a Dice loss for improved overlap optimisation. All models are trained for a uniform budget of 100 epochs with a differential learning rate schedule: the encoder is trained at $6 \times 10^{-5}$ and task heads at $6 \times 10^{-4}$. We use GroupNorm~\cite{wu2018groupnorm} throughout, as the small batch size of 4 makes batch normalisation statistics unreliable. Mixed precision training is employed to reduce memory consumption.

\subsection{Progressive Auxiliary Supervision Ladder}
\label{sec:aux_ladder}

We construct a progressive baseline ladder by incrementally adding auxiliary morphological supervision to the B0 baseline. Each configuration adds exactly one supervision type, with all other training details held constant, ensuring a fair ablation.

\paragraph{B1a: clDice loss.}
B1a adds a soft centerline Dice (clDice) loss~\cite{shit2021cldice} applied to the crack class. The clDice loss is defined as the harmonic mean of soft skeleton precision and soft skeleton recall, computed on differentiable soft skeletons of the prediction and ground truth. To avoid destabilising early training, the clDice loss is activated at epoch 40 with a 5-epoch linear ramp-up, and is weighted at 0.15 relative to the base segmentation losses. This configuration adds no additional parameters, as clDice operates only on the existing segmentation output.

\paragraph{B2: skeleton distance transform regression.}
B2 adds a skeleton distance transform (DT) regression head that predicts the normalised Euclidean distance of each pixel to the nearest crack boundary. The DT target is computed from the ground-truth crack mask using the standard Euclidean distance transform and normalised to $[0, 1]$. The regression head is trained with mean squared error (MSE) loss on the full image (unmasked), with a loss weight of 10.0. This weight was selected through a sweep over values from 1.0 to 20.0 (see Section~\ref{sec:results_aux}). The DT head adds approximately 0.1M parameters to the model.

\paragraph{B3: endpoint and junction detection.}
B3 adds endpoint and junction detection heads that predict Gaussian heatmaps centred on crack keypoints, following a CenterNet-style~\cite{zhou2019centernet} formulation. Endpoints are defined as degree-1 pixels on the crack skeleton, and junctions as degree$\geq$3 pixels. Due to extreme class imbalance (keypoints occupy $<$0.01\% of pixels), each head is trained with a combination of binary cross-entropy with high positive weight and a keypoint Dice loss. B3 adds approximately 0.4M parameters.

\paragraph{B5: full joint supervision.}
B5 combines all auxiliary heads---skeleton DT regression, endpoint and junction detection, and crack width regression---into a single multi-task model. The width head predicts per-pixel crack width using a Softplus activation for non-negative output, trained with MSE loss. Auxiliary losses are activated progressively during training to prevent early-stage interference. B5 adds approximately 0.6M parameters over B0.

\subsection{Post-Hoc Graph Extraction Methods}
\label{sec:graph_extraction}

We evaluate three methods for extracting crack graph representations from model predictions, without any task-specific training for graph prediction.

\paragraph{Method A: skeleton extraction.}
The predicted crack mask is binarised and morphologically skeletonised to produce a one-pixel-wide skeleton. Endpoints (degree-1 pixels) and junctions (degree$\geq$3 pixels) are identified from the skeleton topology, and connected skeleton segments between these keypoints form graph edges. This is the most straightforward approach and serves as the primary extraction method.

\paragraph{Method B: DT ridge detection.}
Ridges in the predicted distance transform map are detected by finding pixels that are local maxima in at least one direction perpendicular to the crack orientation. Pixels with DT values above a threshold form a ridge mask, which is skeletonised and converted to a graph as in Method A.

\paragraph{Method C: DT peak detection.}
Local maxima (peaks) in the predicted DT map are identified as candidate nodes. Peaks are connected by proximity: two peaks form an edge if their Euclidean distance is below a threshold and a high-DT-value path exists between them.

\paragraph{Multi-tier evaluation.}
Graph extraction quality is evaluated at three tolerance tiers. At the \emph{strict} tier (@5\,px), a predicted node must lie within 5 pixels of a ground-truth node to count as a match. The \emph{relaxed} (@10\,px) and \emph{lenient} (@15\,px) tiers allow increasingly generous spatial tolerance. Edge F1 is computed at each tier by matching predicted edges to ground-truth edges based on node correspondence. This multi-tier protocol reveals how sensitive graph quality is to spatial precision.

\subsection{Domain Generalisation Methods}
\label{sec:dg_methods}

We evaluate three standard domain generalisation methods in a single-source pseudo-domain setting.

\paragraph{Pseudo-domain construction.}
Since only a single labelled source domain (DamSegment) is available, we construct pseudo-domains by partitioning training images into three difficulty-based tiers---Easy, Medium, and Hard---based on defect area, number of connected components, and segmentation difficulty. These tiers serve as proxy source domains for DG methods that require multiple training domains.

\paragraph{CORAL~\cite{sun2016coral}.}
Deep CORAL aligns second-order feature statistics (covariance matrices) between pseudo-domain feature distributions. The CORAL loss is applied at encoder stage 2, penalising the Frobenius norm of the difference between covariance matrices of features from different pseudo-domains.

\paragraph{DANN~\cite{ganin2016dann}.}
Domain-Adversarial Neural Networks use a gradient reversal layer (GRL) connected to a domain discriminator MLP. The discriminator is trained to predict which pseudo-domain a feature comes from, while the GRL reverses gradients during backpropagation, encouraging the encoder to produce domain-invariant features. The GRL scaling factor $\alpha$ is scheduled from 0 to 1 over the course of training.

\paragraph{MixStyle~\cite{zhou2021mixstyle, zhou2023mixstyle}.}
MixStyle synthesises new feature statistics by mixing the instance normalisation mean and standard deviation of features from different pseudo-domains using a Beta-distributed mixing coefficient. MixStyle modules are applied as forward hooks on encoder stages 0 and 1, operating directly on intermediate feature maps without additional learnable parameters.

\paragraph{Multi-seed protocol.}
All four methods (ERM baseline, CORAL, DANN, MixStyle) are trained with three random seeds (42, 123, 7), yielding 12 independent runs. Statistical significance is assessed via paired $t$-tests across seeds, with significance declared at $p < 0.05$.


% ============================================================
% SECTION 4: EXPERIMENTAL SETUP
% ============================================================

\section{Experimental Setup}
\label{sec:experimental_setup}

\subsection{Datasets}
\label{sec:datasets}

\paragraph{DamSegment.}
The DamSegment dataset contains 1500 concrete dam inspection images at $640 \times 640$ pixel resolution, resized to $512 \times 512$ for training. Ground-truth annotations are provided as RGB-encoded masks with three semantic classes: background, crack, and spalling. The dataset exhibits severe class imbalance: background accounts for 96.6\% of all pixels, crack for 2.2\%, and spalling for 1.2\%. We split the dataset into 1275 training and 225 validation images. Perceptual hashing is used during dataset preparation to detect and remove near-duplicate images, preventing data leakage between splits.

\paragraph{s2ds.}
The s2ds dataset contains 743 images at $512 \times 512$ pixel resolution from different inspection sites. It is used exclusively as an out-of-distribution (OOD) test set for cross-domain evaluation and is never used during training or hyperparameter selection. This dataset represents the realistic deployment scenario where a model trained on one site must generalise to imagery from unseen sites with different sensors, lighting conditions, and surface textures.

\subsection{Evaluation Metrics}
\label{sec:metrics}

\paragraph{Segmentation metrics.}
The primary segmentation metric is mIoU\textsubscript{fg}, defined as the mean Intersection-over-Union computed over the two foreground classes (crack and spalling), excluding the dominant background class. We also report per-class IoU and boundary F1 (BF1) for crack predictions.

\paragraph{Graph metrics.}
Graph extraction quality is assessed through endpoint F1, junction F1, and edge F1 at three tolerance tiers (strict @5\,px, relaxed @10\,px, lenient @15\,px). Additional metrics include graph edit distance (GED), path continuity, and width mean absolute error (MAE). Paired per-image comparisons between models use the Wilcoxon signed-rank test.

\paragraph{Instance metrics.}
Spalling instance segmentation is evaluated using precision, recall, and F1 computed via connected-component extraction and Hungarian matching at an IoU threshold of 0.5. Both micro-averaged (instance-weighted) and macro-averaged (image-weighted) metrics are reported.

\paragraph{Cross-domain metrics.}
The domain gap is defined as the difference between DamSegment mIoU\textsubscript{fg} and s2ds mIoU\textsubscript{fg}, expressed as an absolute difference.

\paragraph{Statistical tests.}
Multi-seed comparisons use paired $t$-tests with significance declared at $p < 0.05$. We also report Two One-Sided Tests (TOST) equivalence tests with a margin of $\pm$1\% mIoU to assess whether methods are provably equivalent. Per-image paired comparisons use the Wilcoxon signed-rank test.

\subsection{Implementation Details}
\label{sec:implementation}

All models use a SegFormer-B2 encoder pretrained on ImageNet-1K~\cite{deng2009imagenet}, fine-tuned end-to-end. The optimiser is AdamW~\cite{loshchilov2019adamw} with weight decay $1 \times 10^{-4}$ and a cosine annealing learning rate schedule with 5-epoch linear warmup. The batch size is 4 with gradient accumulation to simulate larger effective batches where needed. Mixed precision training (automatic mixed precision) is used throughout. All experiments are conducted on a single NVIDIA GPU.

For multi-seed validation, experiments are repeated with seeds 42, 123, and 7, applied consistently to all sources of randomness (weight initialisation, data shuffling, augmentation). The B0, B1a, and B2\_best configurations, as well as all four DG methods, are validated with all three seeds.

Direct graph prediction experiments (P3) use a SharedFPN at $128 \times 128$ resolution with a CenterNet-style~\cite{zhou2019centernet} node heatmap head and a pairwise edge classifier. The overfit validation uses 16 training images for 200 epochs; full training uses the complete 1275-image training set for 100 epochs.


% ============================================================
% SECTION 5: RESULTS AND DISCUSSION
% ============================================================

\section{Results and Discussion}
\label{sec:results}

\subsection{Auxiliary Supervision Does Not Improve Segmentation}
\label{sec:results_aux}

Table~\ref{tab:single_seed_ladder} presents the single-seed baseline ladder results. At seed 42, B2\_best (skeleton DT, weight=10, MSE loss) achieves the highest mIoU\textsubscript{fg} of 0.683, representing an apparent +1.0\% improvement over B0 (0.673). B1a (clDice) appears to decrease performance to 0.657 ($-$1.6\%), while more complex configurations B3 ($-$0.5\%) and B5 ($-$2.7\%) show degradation relative to B0.

\begin{table}[t]
\centering
\caption{Single-seed baseline ladder results (seed=42). mIoU\textsubscript{fg} is the mean IoU over foreground classes (crack, spalling). All models trained for 100 epochs on DamSegment.}
\label{tab:single_seed_ladder}
\begin{tabular}{llccc}
\toprule
Model & Description & mIoU\textsubscript{fg} & BF1\textsubscript{crack} & Params \\
\midrule
B0 & Mask-only (CE+Dice) & 0.673 & 0.609 & 27.1M \\
B1a & B0 + clDice & 0.657 & --- & 27.1M \\
B2\_best & B0 + skeleton DT (MSE, $w{=}10$) & 0.683 & 0.618 & 27.2M \\
B3 & B0 + endpoint/junction heads & 0.668 & --- & 27.5M \\
B5 & B0 + full joint (skel+kp+width) & 0.646 & --- & 27.7M \\
\bottomrule
\end{tabular}
\end{table}

However, these single-seed results are unreliable. Table~\ref{tab:multiseed_validation} presents the multi-seed validation results for B0, B1a, and B2\_best across three seeds. The multi-seed mean for B2\_best is $0.676 \pm 0.007$, representing only a +0.5\% improvement over B0's $0.671 \pm 0.008$. A paired $t$-test yields $p = 0.145$, which is not significant at the 0.05 level. Similarly, B1a achieves $0.673 \pm 0.015$ ($p = 0.822$), showing no significant difference from B0. The TOST equivalence test also fails to establish equivalence at the $\pm$1\% margin for either method, though this is likely due to the small sample size ($n = 3$) rather than a genuine effect. This comparison is illustrated in Figure~\ref{fig:fig_multiseed_baseline_comparison}.

\begin{table}[t]
\centering
\caption{Multi-seed validation of B0, B1a, and B2\_best on mIoU\textsubscript{fg} (3 seeds). Neither auxiliary method achieves statistically significant improvement over the mask-only baseline.}
\label{tab:multiseed_validation}
\begin{tabular}{lcccc}
\toprule
 & Seed 7 & Seed 42 & Seed 123 & Mean $\pm$ Std \\
\midrule
B0 & 0.681 & 0.663 & 0.669 & $0.671 \pm 0.008$ \\
B1a (clDice) & 0.693 & 0.667 & 0.659 & $0.673 \pm 0.015$ \\
B2\_best (skelDT) & 0.687 & 0.672 & 0.670 & $0.676 \pm 0.007$ \\
\midrule
\multicolumn{5}{l}{\textit{Statistical tests vs B0:}} \\
\midrule
 & $\Delta$ Mean & Paired $t$ $p$ & TOST $p$ & Conclusion \\
\midrule
B1a & $+0.002 \pm 0.009$ & 0.822 (ns) & 0.159 & Not significant \\
B2\_best & $+0.005 \pm 0.003$ & 0.145 (ns) & 0.083 & Not significant \\
\bottomrule
\end{tabular}
\end{table}

\begin{figure}[t]
\centering
\includegraphics[width=\columnwidth]{figures/fig_multiseed_baseline_comparison.png}
\caption{Multi-seed comparison of baseline methods on mIoU\_fg (3 seeds: 42, 123, 7). Error bars show standard deviation; black dots show individual seed values. Neither B1a (clDice, p=0.822) nor B2\_best (skeleton DT, p=0.145) achieves statistically significant improvement over the B0 mask-only baseline. The dashed line indicates the B0 mean for reference.}
\label{fig:fig_multiseed_baseline_comparison}
\end{figure}

Notably, B1a exhibits the highest variance across seeds (std = 0.015, compared to 0.008 for B0), suggesting that the clDice loss makes training less stable. The B2 skeleton DT weight sweep (Table~\ref{tab:dt_weight_sweep}) shows an apparent optimum at weight 10.0 in a single seed, but this advantage does not survive multi-seed validation. The most complex configuration, B5, which stacks all auxiliary heads, performs worst at 0.646, indicating that the additional supervision signals interfere with rather than enhance segmentation learning.

\begin{table}[t]
\centering
\caption{Skeleton DT loss weight sweep (B2 variants, seed=42). The apparent optimum at $w{=}10$ does not survive multi-seed validation.}
\label{tab:dt_weight_sweep}
\begin{tabular}{ccc}
\toprule
Weight & Loss & mIoU\textsubscript{fg} \\
\midrule
1.0 & BCE & 0.664 \\
5.0 & MSE & 0.673 \\
8.0 & MSE & 0.678 \\
10.0 & MSE & 0.683 \\
12.0 & MSE & 0.682 \\
15.0 & MSE & 0.674 \\
20.0 & MSE & 0.669 \\
\bottomrule
\end{tabular}
\end{table}


\subsection{Post-Hoc Graph Extraction Remains Challenging}
\label{sec:results_graph}

Table~\ref{tab:graph_extraction} presents the graph extraction quality for Method A (skeleton extraction) applied to B2\_best predictions, evaluated on 219 crack-containing validation images. Even at the most lenient tolerance tier (@15\,px), edge F1 reaches only 0.277, and the median edge F1 across all images is 0. This indicates that for the majority of images, post-hoc graph extraction fails to produce edges that match the ground-truth graph structure. The graph extraction quality at all three tiers is visualised in Figure~\ref{fig:fig_graph_extraction_edge_f1}.

\begin{table}[t]
\centering
\caption{Post-hoc graph extraction quality using Method A (skeleton extraction) on B2\_best predictions. Mean values over 219 crack-containing images. Even at lenient tolerance, edge F1 remains below 0.28.}
\label{tab:graph_extraction}
\begin{tabular}{lccc}
\toprule
Metric & Strict @5\,px & Relaxed @10\,px & Lenient @15\,px \\
\midrule
Endpoint F1 & 0.290 & 0.433 & 0.509 \\
Junction F1 & 0.392 & 0.405 & 0.408 \\
Edge F1 & 0.092 & 0.220 & 0.277 \\
\midrule
Path continuity & \multicolumn{3}{c}{0.151} \\
GED & \multicolumn{3}{c}{0.379} \\
Width MAE (px) & \multicolumn{3}{c}{2.51} \\
\bottomrule
\end{tabular}
\end{table}

\begin{figure}[t]
\centering
\includegraphics[width=\columnwidth]{figures/fig_graph_extraction_edge_f1.png}
\caption{Post-hoc graph extraction quality measured by edge F1 at three tolerance tiers for the B2\_best model. Method A (skeleton extraction) consistently outperforms DT-based methods B (ridge detection) and C (peak detection). Even at the most lenient tolerance (15 px), edge F1 remains below 0.28, indicating fundamental limitations of post-hoc graph extraction from predicted masks.}
\label{fig:fig_graph_extraction_edge_f1}
\end{figure}

Table~\ref{tab:b0_vs_b2_graph} compares graph extraction quality between B0 and B2\_best using Method A. All Wilcoxon signed-rank $p$-values exceed 0.23, confirming that skeleton DT supervision does not improve post-hoc graph extraction quality. This is a notable negative finding: the DT head was designed to produce features conducive to skeleton extraction, yet it provides no measurable benefit for graph recovery.

\begin{table}[t]
\centering
\caption{Graph extraction comparison: B0 vs B2\_best using Method A. All Wilcoxon signed-rank tests are non-significant ($p > 0.23$), indicating skeleton supervision does not improve graph extraction.}
\label{tab:b0_vs_b2_graph}
\begin{tabular}{lccc}
\toprule
Metric & B0 & B2\_best & Wilcoxon $p$ \\
\midrule
Endpoint F1 @5\,px & 0.283 & 0.290 & 0.720 \\
Edge F1 @5\,px & 0.104 & 0.092 & 0.367 \\
Edge F1 @10\,px & 0.219 & 0.220 & 0.632 \\
Edge F1 @15\,px & 0.268 & 0.277 & 0.234 \\
Edge F1 soft & 0.247 & 0.240 & 0.469 \\
Path continuity & 0.170 & 0.151 & 0.279 \\
\bottomrule
\end{tabular}
\end{table}

Among the three extraction methods applied to B2\_best, Method A (skeleton extraction) consistently outperforms the DT-based alternatives. Method B (DT ridge detection) achieves edge F1 of only 0.053 at strict tolerance and 0.104 at relaxed tolerance. Method C (DT peak detection) performs worst, with edge F1 of 0.003 at strict tolerance. These results indicate that traditional morphological skeletonisation, despite its simplicity, remains the most effective post-hoc graph extraction strategy.

As a complementary evaluation, we assess spalling instance segmentation using connected-component extraction from B2 semantic predictions (Table~\ref{tab:spalling_instance}). This post-hoc approach achieves a micro F1 of 0.567 and macro F1 of 0.482 at IoU threshold 0.5, providing a reasonable baseline without any learned instance segmentation head.

\begin{table}[t]
\centering
\caption{Spalling instance segmentation via post-hoc connected-component extraction from B2 semantic predictions. IoU threshold = 0.5.}
\label{tab:spalling_instance}
\begin{tabular}{lc}
\toprule
Metric & Value \\
\midrule
Micro Precision & 0.563 \\
Micro Recall & 0.571 \\
Micro F1 & 0.567 \\
Macro Precision & 0.495 \\
Macro Recall & 0.484 \\
Macro F1 & 0.482 \\
Mean matched IoU & 0.431 \\
\bottomrule
\end{tabular}
\end{table}


\subsection{Direct Graph Prediction Hits the Auto-Label Ceiling}
\label{sec:results_p3}

We evaluate direct neural graph prediction (P3) using a CenterNet-style~\cite{zhou2019centernet} node heatmap head and pairwise edge classifier, with pre-registered go/no-go gates at each stage. Results are summarised in Table~\ref{tab:p3_overfit} (overfit setting) and Table~\ref{tab:p3_full} (full training), with the comparison illustrated in Figure~\ref{fig:fig_p3_overfit_vs_full}.

\begin{table}[t]
\centering
\caption{P3 overfit results (16 images, 200 epochs). All go/no-go gates pass, confirming the architecture is capable of learning graph structure when given sufficient capacity per sample.}
\label{tab:p3_overfit}
\begin{tabular}{llccc}
\toprule
Gate & Metric & P3a & P3b & Threshold \\
\midrule
1 & Node F1 @5\,px & 0.934 & --- & $\geq 0.90$ \textsc{pass} \\
2 & Edge\textsubscript{gt} F1 @10\,px & --- & 0.917 & $\geq 0.90$ \textsc{pass} \\
3 & Edge F1 @10\,px (pred nodes) & --- & 0.865 & $\geq 0.70$ \textsc{pass} \\
\bottomrule
\end{tabular}
\end{table}

\begin{table}[t]
\centering
\caption{P3 full training results (1275 images, 100 epochs). Node localisation and edge prediction gates fail, revealing the auto-label quality ceiling.}
\label{tab:p3_full}
\begin{tabular}{llcc}
\toprule
Gate & Metric & P3b & Threshold \\
\midrule
4 & Node F1 @5\,px & 0.30 & $>$ post-hoc \textsc{fail} \\
5 & Edge F1 @10\,px & 0.195 & $> 0.30$ \textsc{fail} \\
6 & mIoU\textsubscript{fg} & 0.676 & $\geq 0.673$ \textsc{pass} \\
\bottomrule
\end{tabular}
\end{table}

\begin{figure}[t]
\centering
\includegraphics[width=\columnwidth]{figures/fig_p3_overfit_vs_full.png}
\caption{Comparison of direct graph prediction (P3) performance between overfit setting (16 training images, 200 epochs) and full training (1275 images, 100 epochs). Dotted lines show go/no-go gate thresholds. All gates pass convincingly in the overfit setting, confirming architectural soundness. However, full training fails on node F1 and edge F1 gates, revealing the auto-label quality ceiling: the model cannot generalize when trained on noisy auto-generated graph labels.}
\label{fig:fig_p3_overfit_vs_full}
\end{figure}

In the overfit setting (16 images, 200 epochs), all gates pass convincingly: node F1 reaches 0.934, edge\textsubscript{gt} F1 (using ground-truth nodes) reaches 0.917, and edge F1 from predicted nodes reaches 0.865. This demonstrates that the architecture is fundamentally capable of learning graph structure. However, under full training (1275 images, 100 epochs), performance drops dramatically: node F1 falls to 0.30 and edge F1 to 0.195, both failing their go/no-go gates. Notably, the segmentation gate (mIoU\textsubscript{fg} $\geq$ 0.673) passes at 0.676, confirming that the graph prediction failure is not caused by degraded segmentation.

A hybrid experiment (P3h) provides further diagnostic insight. When skeleton-extracted nodes are used in place of predicted nodes and only the edge classifier operates on learned features, edge\textsubscript{gt} F1 reaches 0.654---substantially higher than P3b's 0.454. This confirms that the edge classifier generalises well and that node localisation is the primary bottleneck. The root cause is the quality of automatically generated training labels: the mask-to-graph pipeline used to produce node and edge labels from segmentation masks introduces systematic noise that limits what the model can learn at scale. Gold-standard graph annotations are needed to advance beyond this ceiling.


\subsection{Domain Generalisation Methods Are Ineffective Under Pseudo-Domains}
\label{sec:results_dg}

Table~\ref{tab:dg_multiseed} presents the multi-seed domain generalisation results. All four methods achieve similar in-domain performance on DamSegment (mIoU\textsubscript{fg} ranging from $0.668 \pm 0.010$ for MixStyle to $0.675 \pm 0.007$ for ERM) and similar out-of-distribution performance on s2ds (mIoU\textsubscript{fg} ranging from $0.115 \pm 0.012$ for ERM to $0.120 \pm 0.011$ for DANN). The domain gap remains stable at approximately 55\% across all methods. A paired $t$-test comparing MixStyle to ERM on s2ds mIoU\textsubscript{fg} yields $p = 0.878$, confirming statistical equivalence. The multi-seed results are visualised in Figure~\ref{fig:fig_dg_multiseed_domain_gap}.

\begin{table}[t]
\centering
\caption{Domain generalisation results (3 seeds: 42, 123, 7). All methods are statistically equivalent to ERM ($p = 0.878$ for MixStyle vs ERM). The domain gap of ${\sim}$55\% is not reduced by any DG method.}
\label{tab:dg_multiseed}
\begin{tabular}{llccc}
\toprule
Method & Type & DamSeg mIoU\textsubscript{fg} & s2ds mIoU\textsubscript{fg} & Domain Gap \\
\midrule
ERM & baseline & $0.675 \pm 0.007$ & $0.115 \pm 0.012$ & $0.560 \pm 0.015$ \\
CORAL & feature align & $0.672 \pm 0.006$ & $0.118 \pm 0.010$ & $0.554 \pm 0.013$ \\
DANN & adversarial & $0.674 \pm 0.009$ & $0.120 \pm 0.011$ & $0.553 \pm 0.015$ \\
MixStyle & style mixing & $0.668 \pm 0.010$ & $0.116 \pm 0.006$ & $0.552 \pm 0.016$ \\
\bottomrule
\end{tabular}
\end{table}

\begin{figure}[t]
\centering
\includegraphics[width=\columnwidth]{figures/fig_dg_multiseed_domain_gap.png}
\caption{Domain generalization results across 3 seeds for four methods on in-domain (DamSegment) and out-of-domain (s2ds) evaluation. All methods achieve similar in-domain performance (${\sim}$0.67 mIoU\_fg) and similar OOD performance (${\sim}$0.12 mIoU\_fg), with a stable domain gap of approximately 55\%. No DG method significantly reduces this gap (p=0.878 for MixStyle vs ERM). Dashed lines indicate ERM baseline means.}
\label{fig:fig_dg_multiseed_domain_gap}
\end{figure}

The single-seed comparison at seed 42 (Table~\ref{tab:dg_singleseed}) paints a strikingly different picture: MixStyle appears to improve s2ds performance by +52\% (0.167 vs 0.110), while CORAL appears to gain +21\% (0.133 vs 0.110). These single-seed results would support a positive narrative about DG methods for cross-site crack segmentation. The multi-seed validation completely overturns this narrative, revealing that the apparent improvements were artefacts of seed-dependent variance.

\begin{table}[t]
\centering
\caption{Single-seed DG comparison (seed=42 only) that would suggest DG methods help. Multi-seed validation (Table~\ref{tab:dg_multiseed}) overturns these conclusions.}
\label{tab:dg_singleseed}
\begin{tabular}{lccc}
\toprule
Method & DamSeg mIoU\textsubscript{fg} & s2ds mIoU\textsubscript{fg} & Domain Gap \\
\midrule
ERM & 0.664 & 0.110 & 0.554 \\
CORAL & 0.666 & 0.133 & 0.533 \\
DANN & 0.666 & 0.121 & 0.545 \\
MixStyle & 0.658 & 0.167 & 0.490 \\
\bottomrule
\end{tabular}
\end{table}

The failure of DG methods in this setting is attributable to the pseudo-domain construction strategy. Difficulty-based data tiers (Easy/Medium/Hard) capture variation in defect complexity but not genuine distributional shifts in imaging conditions, surface textures, or sensor characteristics. DG methods are designed to learn invariance across domains with distinct visual characteristics; when pseudo-domains share the same underlying distribution, the domain-invariance objective provides no useful learning signal. This suggests that bridging real cross-site gaps requires genuine multi-source data collection from diverse inspection sites rather than algorithmic solutions applied to single-source data.


\subsection{The Danger of Single-Seed Evaluation}
\label{sec:results_singleseed}

The preceding sections document three cases where single-seed evaluation produced conclusions that were overturned by multi-seed validation. We summarise these cases and discuss their implications, with all three illustrated in Figure~\ref{fig:fig_single_vs_multiseed_overturned}.

\begin{figure}[t]
\centering
\includegraphics[width=\columnwidth]{figures/fig_single_vs_multiseed_overturned.png}
\caption{Three cases where single-seed conclusions were overturned by multi-seed validation. (a) B2 skeleton DT appeared to improve mIoU by +1.0\% on seed 42, but multi-seed reveals +0.5\% +/- 0.3\% (p=0.145, not significant). (b) B1a clDice appeared to harm mIoU by -1.6\%, but multi-seed shows +0.2\% +/- 0.9\% (p=0.822, not significant). (c) MixStyle appeared to improve OOD performance by +52\%, but multi-seed shows approximately +1\% (p=0.878, not significant).}
\label{fig:fig_single_vs_multiseed_overturned}
\end{figure}

\paragraph{Case 1: B2 skeleton DT.}
A single seed (42) showed B2\_best outperforming B0 by +1.0\% mIoU\textsubscript{fg} (0.683 vs 0.673). This magnitude is typical of claimed improvements in the crack segmentation literature and would normally be considered a positive result. Multi-seed evaluation reduced the effect to +0.5\% $\pm$ 0.3\% ($p = 0.145$), revealing that approximately half of the apparent improvement was attributable to seed variance.

\paragraph{Case 2: B1a clDice.}
A single seed (42) showed clDice degrading performance by $-$1.6\% (0.657 vs 0.673), which would lead to the conclusion that clDice harms crack segmentation. Multi-seed evaluation reversed the direction entirely, showing +0.2\% $\pm$ 0.9\% ($p = 0.822$). The high variance (std = 0.015 vs B0's 0.008) explains why a single seed could produce such a misleading result.

\paragraph{Case 3: MixStyle OOD.}
A single seed (42) showed MixStyle improving s2ds performance from 0.110 to 0.167, an apparent +52\% relative improvement. Multi-seed evaluation reduced this to a negligible +1\% ($p = 0.878$), demonstrating that the apparent DG benefit was entirely an artefact of seed-dependent variance on the out-of-distribution test set.

These three cases share a common pattern: the effect sizes attributable to seed variance (0.5--5.7\% absolute, up to 52\% relative) are comparable to or larger than the improvements typically reported in the literature. This is consistent with findings by Picard~\cite{picard2021seed}, who showed that seed selection can produce performance differences of up to 1.8\% on ImageNet, and by Henderson et al.~\cite{henderson2018rl}, who documented similar fragility in reinforcement learning benchmarks. In the concrete defect detection domain, Alsabbagh and Almasri~\cite{alsabbagh2024randomness} recently demonstrated that randomness-induced variance can rival reported improvements in civil engineering and medical imaging applications.

The implications for the broader crack segmentation literature are concerning. If single-seed evaluation can produce misleading results of the magnitude documented here, then a non-trivial fraction of published claims of improvement may not survive multi-seed validation. We argue that multi-seed validation with at least three seeds, accompanied by appropriate statistical tests, should be adopted as a minimum standard for claiming improvement in this domain.


% ============================================================
% SECTION 6: CONCLUSIONS
% ============================================================

\section{Conclusions}
\label{sec:conclusions}

This study presented a systematic, multi-seed evaluation of auxiliary morphological supervision for crack segmentation, post-hoc graph extraction, direct neural graph prediction, and standard domain generalisation methods for cross-site concrete defect detection. The central finding is negative on all fronts:

\begin{enumerate}
    \item \textbf{Auxiliary morphological supervision does not improve segmentation.} Neither clDice (+0.2\% $\pm$ 0.9\%, $p = 0.822$) nor skeleton distance transform regression (+0.5\% $\pm$ 0.3\%, $p = 0.145$) achieves statistically significant improvement over the mask-only baseline under multi-seed validation. More complex multi-head configurations (B3, B5) degrade performance.

    \item \textbf{Post-hoc graph extraction remains fundamentally challenging.} Median edge F1 is 0 across all extraction methods, and skeleton DT supervision does not improve graph extraction quality (Wilcoxon $p > 0.23$ for all metrics). Morphological skeletonisation (Method A) outperforms DT-based alternatives (Methods B, C).

    \item \textbf{Direct neural graph prediction is limited by auto-label quality.} The architecture is capable of learning graph structure (overfit gates pass at node F1 = 0.934, edge F1 = 0.865), but auto-generated training labels create a quality ceiling that prevents generalisation (full-train node F1 = 0.30, edge F1 = 0.195). Gold-standard graph annotations are needed.

    \item \textbf{Standard DG methods are ineffective under pseudo-domain single-source training.} ERM, CORAL, DANN, and MixStyle are statistically equivalent ($p = 0.878$), with a stable domain gap of approximately 55\%. Difficulty-based pseudo-domains do not capture genuine cross-site distribution shifts.
\end{enumerate}

The most impactful contribution of this work is methodological: we documented three concrete cases where single-seed evaluation produced misleading conclusions---a positive result that was halved, a negative result whose direction reversed, and a +52\% relative improvement that vanished entirely. These findings demonstrate that the effect sizes typically reported in the crack segmentation literature are within the range of seed-dependent variance, calling into question the reliability of single-seed claims of improvement.

Based on these findings, we offer two practical recommendations. First, for concrete defect segmentation, the mask-only baseline (B0) should be the default choice: it is the simplest configuration, achieves equivalent performance to all auxiliary supervision variants, and avoids the added complexity and instability of multi-head training. Second, for cross-site generalisation, investment in multi-site labelled data collection is more likely to yield progress than the application of DG methods to single-source data with constructed pseudo-domains.

Future work should pursue three directions: (1) gold-standard graph annotations for a representative subset of crack images, enabling direct neural graph prediction beyond the auto-label ceiling; (2) true multi-source domain generalisation with labelled data from multiple inspection sites, sensors, and environmental conditions; and (3) learned instance segmentation heads for spalling, building on the reasonable post-hoc connected-component baseline (micro F1 = 0.567) established in this study.


% ============================================================
% DATA AVAILABILITY
% ============================================================

\section*{Data Availability}

The DamSegment dataset used in this study is publicly available. The s2ds dataset is used under licence for research purposes. Code and trained model weights will be released upon acceptance.


% ============================================================
% BIBLIOGRAPHY
% ============================================================

\bibliographystyle{elsarticle-num}
\bibliography{refs}

\end{document}
